% Generado por generar_tablas.py -- NO editar a mano.

\begin{table}[htbp]
\centering
\caption{Comparación de las tres familias de mitigación sobre el conjunto de
prueba. Menor $\Delta$TPR indica mayor equidad entre entidades.}
\label{tab:mitigacion}
\small\setlength{\tabcolsep}{5pt}
\begin{tabular}{l r r r r}
\toprule
Intervención & $\Delta$TPR & $\Delta$FPR & Selec. & Exact.\ bal. \\
\midrule
Línea de base (umbral global) & 0.294 & 0.087 & 0.100 & 0.868 \\
Pre: reponderación por entidad & 0.290 & 0.082 & 0.099 & 0.865 \\
In: \emph{equalized odds} (Fairlearn)\textsuperscript{a} & 0.445 & -- & -- & -- \\
Post: umbral por entidad\textsuperscript{b} & 0.140 & 0.101 & 0.102 & 0.871 \\
Post: umbral por entidad (en muestra)\textsuperscript{c} & 0.022 & 0.107 & 0.101 & 0.869 \\
\bottomrule
\end{tabular}

\vspace{0.4em}
\begin{minipage}{0.95\linewidth}\footnotesize
\textsuperscript{a} Evaluado en una submuestra del bloque de entrenamiento
($n=\numprint{120000}$) y con un clasificador lineal; no es comparable uno a uno
con las demás filas. La reducción devuelve una distribución sobre clasificadores:
sobre 20 realizaciones del mismo modelo
ajustado, la brecha tiene mediana 0.397 y rango
[0.245, 0.628], intervalo
que contiene a la línea de base.\\
\textsuperscript{b} Umbrales estimados en el bloque de calibración y evaluados
en el de prueba: es la cifra desplegable.\\
\textsuperscript{c} Umbrales estimados sobre el propio conjunto de prueba:
cota optimista, reportada solo para dimensionar el sesgo de ajuste en muestra.
\end{minipage}
\end{table}
